\vspace{-8pt}
\section{Problem Formulation}
\vspace{-3pt}
\label{sec:background}
We first motivate a new abstraction, the \emph{\process~(\proc)}, as an alternative to the Markov Decision Process~(MDP) commonly used in RL, and then show an instantiation of a \proc with diffusion models.

\vspace{-5pt}
\subsection{Markov Decision Process}\label{subsec:mdp}
\vspace{-3pt}

The Markov Decision Process~\citep{puterman1994markov} is a broad abstraction used to formulate many sequential decision making problems.
Many RL algorithms have been derived from MDPs with empirical success~\citep{haarnoja2018soft,hafner2020mastering,zhang2022making}, but existing algorithms are typically unable to combinatorially generalize across different environments. Such difficulty %
can be traced back to certain aspects of the underlying MDP abstraction: %

\begin{itemize}[leftmargin=2em, nosep]
    \item[{\bf i)}] The lack of a universal state interface across different control environments. In fact, since different environments typically have separate underlying state spaces, one would need to a construct a complex state representation to represent all environments, making learning difficult.  
    \item[{\bf ii)}] The explicit requirement of a real-valued reward function in an MDP. The RL problem is usually defined as maximizing the accumulated reward in an MDP. However, in many practical applications, how to design and transfer rewards is unclear and different across environments.   
    \item[{\bf iii)}] The dynamics model in an MDP is environment and agent dependent. Specifically, the dynamics model $T(s'|s, a)$ that characterizes the transition between states $(s, s')$ under action $a$ is explicitly dependent to the environment and action space of the agent, which can be significantly different between different agents and tasks.
\end{itemize}

\vspace{-5pt}
\subsection{\process}\label{subsec:udp}
\vspace{-3pt}

These difficulties inspire us to construct an alternative abstraction for unified sequential decision making across diverse environments. Our abstraction, termed \emph{\process~(\proc)}, exploits images as a universal interface across environments, texts as task specifiers to avoid reward design, and a task-agnostic planning module separated from environment-dependent control to enable knowledge sharing and generalization. %

Formally, we define a \proc to be a tuple $\Gcal = \langle \Xcal, \Ccal, H, \rho\rangle$, %
where $\Xcal$ denotes the observation space of images, $\Ccal$ denotes the space of textual task descriptions, $H\in \Ncal$ is a finite horizon length, and $\rho(\cdot| x_0, c): \Xcal \times\Ccal\rightarrow \Delta(\Xcal^H)$ 
is a conditional video generator.
That is, $\rho(\cdot|x_o,c)\in\Delta(\Xcal^H)$ is a conditional distribution over $H$-step image sequences determined by the first frame $x_0$ and the task description $c$. 
Intuitively, $\rho$ synthesizes $H$-step image trajectories that illustrate possible paths for completing a target task $c$. For simplicity, we focus on finite horizon, episodic tasks.

Given a \proc $\Gcal$, 
we define a trajectory-task conditioned policy $\pi(\cdot| \{x_h\}_{h=0}^H, c): \Xcal^{H+1}\times\Ccal\rightarrow \Delta(\Acal^{H})$ to be a conditional distribution over $H$-step action sequences $\Acal^H$.
Ideally, $\pi(\cdot| \{x_h\}_{h=0}^H, c)$ specifies a conditional distribution of action sequences that achieves the given trajectory $\{x_h\}_{h=0}^H$ in the \proc $\Gcal$ for the given task $c$.
To achieve such an alignment, we will consider an offline RL scenario where we have access to a dataset of existing experience $\Dset = \{(x_i, a_i)_{i=0}^{H-1}, x_H, c\}_{j=1}^n$
from which both $\rho(\cdot|x_0, c)$ and $\pi(\cdot| \{x_h\}_{h=0}^H, c)$ can be estimated.


\looseness=-1
In contrast to an MDP, a \proc directly models video-based trajectories and bypasses the need to specify a reward function beyond the textual task description. Since the space of video observations $\Xcal^H$ and task descriptions $\Ccal$ are both naturally shared across environments and easily interpretable by humans, any video-based planner $\rho(\cdot|x_0, c)$ can be conveniently \emph{reused}, \emph{transferred} and \emph{debugged}. 
Another benefit of a \proc over an MDP is that \proc isolates the video-based planning with $\rho(\cdot|x_0, c)$ from the deferred action selection using $\pi(\cdot|\{x_h\}_{h=0}^H, c)$. This design choice isolates planning decisions from action-specific mechanisms, allowing the planner to be environment and agent agnostic. %

\proc can be understood as implicitly planning over an MDP and directly outputting an optimal trajectory based on the given instructions. 
Such a \proc abstraction bypasses reward design, state extraction and explicit planning, and allows for \emph{non-Markovian} modeling of an image-based state space.
However, learning a planner in \proc requires videos and task descriptions, whereas traditional MDPs do not require such data, so whether an MDP or \proc is more suitable for a given task depends on what types of training data are available. 
Although the non-Markovian model and the requirement of video and text data induce additional difficulties in \proc comparing to MDP, it is possible to leverage existing large text-video models that have been pretrained on massive, web-scale datasets to alleviate these complexities. 



\vspace{-3pt}
\subsection{Diffusion Models for \proc}\label{subsec:diffusion} 
\vspace{-3pt}

Let $\tau = [x_1, \ldots, x_H]\in \Xcal^{H}$ denote a sequence of images. 
We leverage the significant recent advances in diffusion models for capturing the conditional distribution $\rho(\tau |x_0, c)$, which we will leverage as a text and initial-frame conditioned video generator in a \proc. 
We emphasize that the \proc formulation is also compatible with other probabilistic models, such as a variational autoencoder~\citep{kingma2013vae}, energy-based model~\citep{du2019implicit,dai2019exponential}, or generative adversarial network~\citep{goodfellow2014gan}. 
For completeness we briefly cover the core formulation at a high-level, but defer details to background references~\citep{sohldickstein2015nonequilibrium}. 

We start with an unconditional model. A continuous-time diffusion model defines a forward process $q_k(\tau_k|\tau) = \mathcal{N}(\cdot; \alpha_k\tau, \sigma_k^2 I)$, where $k \in [0, 1]$ and $\alpha_k, \sigma_k^2$ are scalars with predefined schedules. A generative process $p(\tau)$ is also defined, which reverses the forward process by learning a denoising model $s(\tau_k, k)$. Correspondingly $\tau$ can be generated by simulating this reverse process with an ancestral sampler~\citep{ho2020denoising} or numerical integration~\citep{song2020denoising}. In our case, the unconditional model needs to be further adapted to condition on both the text instruction $c$ and the initial image $x_0$. Denote the conditional denoiser as $s(\tau_k, k|c, x_0)$. 
We leverage classifier-free guidance~\citep{ho2022classifier} and use $\hat{s}(\tau_k, k|c, x_0) = (1+\omega) s(\tau_k, k|c, x_0) - \omega s(\tau_k, k)$ as the denoiser in the reverse process for sampling, where $\omega$ controls the strength of the text and first-frame conditioning.